\documentclass[a4paper,12pt]{article}

\usepackage{geometry}
\geometry{top=20mm,bottom=20mm,left=25mm,right=25mm}

\usepackage[utf8]{inputenc}
\usepackage[T2A]{fontenc}
\usepackage[russian]{babel}
\usepackage{amsmath}

\usepackage{xcolor}
\usepackage{listings}
\usepackage{hyperref}

\usepackage{booktabs}
\usepackage{array}

\definecolor{codebg}{RGB}{248,248,248}
\definecolor{codeframe}{RGB}{210,210,210}
\definecolor{codekeyword}{RGB}{0,76,153}
\definecolor{codecomment}{RGB}{0,128,0}
\definecolor{codestring}{RGB}{163,21,21}
\definecolor{codenumber}{RGB}{120,120,120}

\lstdefinestyle{common}{
    backgroundcolor=\color{codebg},
    frame=single,
    rulecolor=\color{codeframe},
    basicstyle=\ttfamily\small,
    keywordstyle=\color{codekeyword}\bfseries,
    commentstyle=\color{codecomment}\itshape,
    stringstyle=\color{codestring},
    numberstyle=\tiny\color{codenumber},
    numbers=left,
    numbersep=10pt,
    breaklines=true,
    showstringspaces=false,
    keepspaces=true,
    columns=fullflexible,
    tabsize=4
}

\lstdefinestyle{cppstyle}{
    style=common,
    language=C++
}

\lstdefinestyle{pythonstyle}{
    style=common,
    language=Python
}

\begin{document}

\begin{titlepage}
    \centering
    {\large МИНОБРНАУКИ РОССИИ\par}
    \vspace{0.5em}
    {\large САНКТ-ПЕТЕРБУРГСКИЙ ГОСУДАРСТВЕННЫЙ\par}
    {\large ЭЛЕКТРОТЕХНИЧЕСКИЙ УНИВЕРСИТЕТ\par}
    {\large «ЛЭТИ» ИМ. В.И. УЛЬЯНОВА (ЛЕНИНА)\par}
    \vspace{1em}
    {\large Кафедра Информационных систем\par}

    \vspace{6em}

    {\LARGE \textbf{Практическая работа}\par}
    \vspace{0.8em}
    {\large по дисциплине «Теория принятия решений»\par}
    \vspace{0.8em}
    {\large Тема: «Применение методов линейного программирования для решения задачи поиска оптимального плана заказа стали на завод»\par}

    \vfill

    \begin{flushleft}
        Студент гр. 3376 \hfill Михайлов Н.Ф.\\[0.5em]
        Преподаватель   \hfill Пономарёв А.В.
    \end{flushleft}

    \vspace{3em}

    Санкт-Петербург\par
    2026
\end{titlepage}

\tableofcontents
\newpage

\section{Введение}
Целью данной практической работы является освоение метода линейного
программирования в задаче оптимального распределения стали между производством,
воспроизводством ресурса и расширением производственных мощностей. \\
Для этого необходимо:
\begin{enumerate}
    \item формализовать задачу распределения стали по годам;
    \item записать переменные, целевую функцию и ограничения;
    \item представить динамику запасов и производственных мощностей в виде
    линейных ограничений;
    \item реализовать численное решение задачи линейного программирования;
    \item исследовать влияние начальных условий на структуру и результат решения.
\end{enumerate}

\section{Постановка задачи}
В пункте $B_2$ сталь поступает на условный производственный комплекс,
состоящий из сталелитейного и станкостроительного заводов. Комплекс
функционирует в течение 5 лет.

Начальный запас стали составляет 3800 т. Исходные производственные мощности:
\begin{itemize}
    \item сталелитейный завод --- 2500 т стали в год;
    \item станкостроительный завод --- 1700 станков в год.
\end{itemize}

Сталь может быть распределена между четырьмя направлениями:
\begin{itemize}
    \item производство станков;
    \item производство стали;
    \item расширение мощности сталелитейного завода;
    \item расширение мощности станкостроительного завода.
\end{itemize}

Для производства одного станка требуется 1.4 т стали. Каждая тонна стали,
направленная на производство стали, обеспечивает выпуск 2.4 т стали.
Каждая тонна стали, направленная на расширение сталелитейного завода,
увеличивает его мощность на 0.1 т в год. Для увеличения мощности
станкостроительного завода на один станок в год требуется 11 т стали.

Решение о распределении стали реализуется в конце очередного года, поэтому
увеличенные производственные мощности используются начиная со следующего года.
Станкостроительный завод не может получать более половины имеющегося запаса
стали.

Требуется:
\begin{enumerate}
    \item определить план распределения стали, обеспечивающий максимальный
    суммарный выпуск станков за 5 лет;
    \item исследовать влияние начального запаса стали в диапазоне 1000--10000 т
    на структуру оптимального решения и суммарный выпуск станков;
    \item исследовать влияние начальной мощности станкостроительного завода
    в диапазоне 100--5000 станков в год на структуру оптимального решения и
    суммарный выпуск станков.
\end{enumerate}

\section{Формализация задачи}

\subsection{Формальная постановка}

\textbf{Переменные.}

Плановый горизонт составляет 5 лет. Для каждого года $i = 1,\dots,5$
вводятся управляющие переменные:
\begin{itemize}
    \item $x_i$ --- количество стали, направляемое на производство станков, т;
    \item $y_i$ --- количество стали, направляемое на производство стали, т;
    \item $z_i$ --- количество стали, направляемое на расширение
    сталелитейного завода, т;
    \item $w_i$ --- количество стали, направляемое на расширение
    станкостроительного завода, т.
\end{itemize}

Также вводятся переменные состояния на начало каждого года:
\begin{itemize}
    \item $s_i$ --- запас стали в начале года $i$, т;
    \item $p_i^{steel}$ --- мощность сталелитейного завода, т стали в год;
    \item $p_i^{mach}$ --- мощность станкостроительного завода, станков в год.
\end{itemize}

Состояния задаются для $i = 1,\dots,6$, так как после пятого года необходимо
описать итоговый запас и итоговые мощности.

Цель задачи --- максимизировать суммарный выпуск станков за 5 лет:
\begin{equation}
\max Z = \sum_{i=1}^{5} \frac{x_i}{1.4}.
\end{equation}

Начальные условия:
\begin{equation}
s_1 = 3800,
\end{equation}
\begin{equation}
p_1^{steel} = 2500,
\end{equation}
\begin{equation}
p_1^{mach} = 1700.
\end{equation}

Для каждого года $i = 1,\dots,5$ выполняются следующие ограничения.

Баланс распределения стали:
\begin{equation}
x_i + y_i + z_i + w_i \le s_i.
\end{equation}

Станкостроительный завод не может получать более половины имеющегося запаса
стали. Так как сталь поступает на него как для выпуска станков, так и для
расширения мощности, ограничение имеет вид:
\begin{equation}
x_i + w_i \le 0.5s_i.
\end{equation}

Выпуск станков ограничен мощностью станкостроительного завода:
\begin{equation}
\frac{x_i}{1.4} \le p_i^{mach}.
\end{equation}

Выпуск стали ограничен мощностью сталелитейного завода:
\begin{equation}
2.4y_i \le p_i^{steel}.
\end{equation}

Динамика запаса стали:
\begin{equation}
s_{i+1} = s_i - x_i - y_i - z_i - w_i + 2.4y_i.
\end{equation}

Или, после приведения подобных слагаемых:
\begin{equation}
s_{i+1} = s_i - x_i + 1.4y_i - z_i - w_i.
\end{equation}

Динамика мощности сталелитейного завода:
\begin{equation}
p_{i+1}^{steel} = p_i^{steel} + 0.1z_i.
\end{equation}

Динамика мощности станкостроительного завода:
\begin{equation}
p_{i+1}^{mach} = p_i^{mach} + \frac{w_i}{11}.
\end{equation}

Все переменные неотрицательны:
\begin{equation}
x_i, y_i, z_i, w_i, s_i, p_i^{steel}, p_i^{mach} \ge 0.
\end{equation}

\subsection{Особенность численного решения}

Несмотря на поэтапную структуру задачи, все ограничения и переходы состояния
являются линейными. Поэтому на практике задача решается как одна задача
линейного программирования, развернутая по всем годам планового периода.

Вектор неизвестных включает управляющие переменные $x_i$, $y_i$, $z_i$, $w_i$
для $i = 1,\dots,5$, а также переменные состояния $s_i$, $p_i^{steel}$,
$p_i^{mach}$ для $i = 1,\dots,6$. После этого целевая функция и все ограничения
записываются в матричной форме, пригодной для передачи в LP-решатель.

Такой подход не требует дискретизации состояний и управлений и позволяет
получить точное решение линейной модели.

\section{Решение подзадач}
\subsection{Подзадача 1. Базовое решение}

Для исходных значений $s_1 = 3800$, $p_1^{steel} = 2500$,
$p_1^{mach} = 1700$ получен следующий оптимальный план.

\begin{center}
\small
\begin{tabular}{ccccccccc}
\toprule
Год & $s_i$ & $p_i^{steel}$ & $p_i^{mach}$ & $x_i$ & $y_i$ & $z_i$ & $w_i$ & Выпуск \\
\midrule
1 & 3800.00 & 2500.00 & 1700.00 & 1900.00 & 1041.67 & 0.00 & 0.00 & 1357.14 \\
2 & 3358.33 & 2500.00 & 1700.00 & 1679.17 & 1041.67 & 0.00 & 0.00 & 1199.40 \\
3 & 3137.50 & 2500.00 & 1700.00 & 1568.75 & 1041.67 & 0.00 & 0.00 & 1120.54 \\
4 & 3027.08 & 2500.00 & 1700.00 & 1513.54 & 1041.67 & 0.00 & 0.00 & 1081.10 \\
5 & 2971.88 & 2500.00 & 1700.00 & 1485.94 & 0.00 & 0.00 & 0.00 & 1061.38 \\
\bottomrule
\end{tabular}
\end{center}

Суммарный выпуск за 5 лет составил:
\[
Z = 5819.57 \text{ станков}.
\]

Итоговое состояние комплекса:
\[
s_6 = 1485.94,\quad p_6^{steel} = 2500,\quad p_6^{mach} = 1700.
\]

\textbf{Вывод.}

В базовом решении расширение мощностей не используется. В первые четыре года
сталелитейный завод работает на пределе мощности:
\[
2.4 \cdot 1041.67 \approx 2500.
\]
Кроме того, в каждом году на производство станков направляется половина
имеющегося запаса стали. Следовательно, основными ограничивающими факторами
являются запас стали и мощность сталелитейного завода, а начальная мощность
станкостроительного завода для базового варианта достаточна.

\subsection{Подзадача 2. Анализ начального запаса стали}

Проводилось изменение начального запаса стали в диапазоне от 1000 до 10000 т.
Остальные начальные параметры оставались неизменными:
\[
p_1^{steel} = 2500,\quad p_1^{mach} = 1700.
\]

\begin{center}
\small
\begin{tabular}{ccccccccc}
\toprule
Нач. запас & Выпуск & $x_1$ & $y_1$ & $z_1$ & $w_1$ & $s_6$ & $p_6^{steel}$ & $p_6^{mach}$ \\
\midrule
1000 & 3820.68 & 0.00 & 1000.00 & 0.00 & 0.00 & 1426.04 & 2500.00 & 1700.00 \\
2000 & 4572.17 & 958.33 & 1041.67 & 0.00 & 0.00 & 1432.29 & 2500.00 & 1700.00 \\
3000 & 5266.00 & 1500.00 & 1041.67 & 0.00 & 0.00 & 1460.94 & 2500.00 & 1700.00 \\
4000 & 5957.96 & 2000.00 & 1041.67 & 0.00 & 0.00 & 1492.19 & 2500.00 & 1700.00 \\
5000 & 6644.57 & 2380.00 & 1041.67 & 0.00 & 0.00 & 1530.94 & 2500.00 & 1700.00 \\
6000 & 7300.00 & 2380.00 & 1041.67 & 0.00 & 0.00 & 1613.33 & 2500.00 & 1700.00 \\
7000 & 7889.58 & 2380.00 & 1041.67 & 0.00 & 0.00 & 1787.92 & 2500.00 & 1700.00 \\
8000 & 8340.48 & 2380.00 & 1041.67 & 0.00 & 0.00 & 2156.67 & 2500.00 & 1700.00 \\
9000 & 8622.96 & 2380.00 & 1041.67 & 0.00 & 338.15 & 2423.04 & 2500.00 & 1730.74 \\
10000 & 8845.19 & 2380.00 & 1041.67 & 0.00 & 949.26 & 2500.81 & 2500.00 & 1786.30 \\
\bottomrule
\end{tabular}
\end{center}

\textbf{Вывод.}

При малом начальном запасе стали оптимальный план направляет значительную часть
ресурса на производство стали. Например, при запасе 1000 т в первый год весь
ресурс направляется на производство стали, а выпуск станков в первый год
отсутствует.

При увеличении начального запаса суммарный выпуск возрастает. Начиная с запаса
5000 т величина $x_1$ достигает 2380 т, что соответствует ограничению мощности
станкостроительного завода:
\[
1700 \cdot 1.4 = 2380.
\]
Дальнейший рост начального запаса уже не может увеличить выпуск первого года
за счет $x_1$, поэтому при больших запасах появляется расширение
станкостроительного завода ($w_1 > 0$).

\subsection{Подзадача 3. Анализ начальной мощности станкостроительного завода}

Проводилось изменение начальной мощности станкостроительного завода в диапазоне
от 100 до 5000 станков в год. Начальный запас стали и мощность сталелитейного
завода оставались равными:
\[
s_1 = 3800,\quad p_1^{steel} = 2500.
\]

\begin{center}
\small
\begin{tabular}{ccccccccc}
\toprule
$p_1^{mach}$ & Выпуск & $x_1$ & $y_1$ & $z_1$ & $w_1$ & $s_6$ & $p_6^{steel}$ & $p_6^{mach}$ \\
\midrule
100 & 1764.58 & 140.00 & 1041.67 & 0.00 & 1760.00 & 1126.26 & 2554.60 & 551.15 \\
600 & 3704.18 & 840.00 & 1041.67 & 0.00 & 1060.00 & 1327.87 & 2521.75 & 841.21 \\
1100 & 5587.41 & 1540.00 & 1041.67 & 0.00 & 240.37 & 1570.59 & 2500.00 & 1121.85 \\
1600 & 5819.57 & 1900.00 & 1041.67 & 0.00 & 0.00 & 1485.94 & 2500.00 & 1600.00 \\
2100 & 5819.57 & 1900.00 & 1041.67 & 0.00 & 0.00 & 1485.94 & 2500.00 & 2100.00 \\
2600 & 5819.57 & 1900.00 & 1041.67 & 0.00 & 0.00 & 1485.94 & 2500.00 & 2600.00 \\
3100 & 5819.57 & 1900.00 & 1041.67 & 0.00 & 0.00 & 1485.94 & 2500.00 & 3100.00 \\
3600 & 5819.57 & 1900.00 & 1041.67 & 0.00 & 0.00 & 1485.94 & 2500.00 & 3600.00 \\
4100 & 5819.57 & 1900.00 & 1041.67 & 0.00 & 0.00 & 1485.94 & 2500.00 & 4100.00 \\
4600 & 5819.57 & 1900.00 & 1041.67 & 0.00 & 0.00 & 1485.94 & 2500.00 & 4600.00 \\
\bottomrule
\end{tabular}
\end{center}

\textbf{Вывод.}

При малой начальной мощности станкостроительного завода существенная часть
стали направляется на его расширение. Так, при мощности 100 станков в год в
первый год на расширение направляется 1760 т стали. Это позволяет увеличить
выпуск в последующие годы.

При росте начальной мощности суммарный выпуск увеличивается до значения
5819.57 станков. Начиная примерно с мощности 1600 станков в год дальнейшее
увеличение начальной мощности уже не повышает результат. Это означает, что
после этого порога ограничивающим фактором становится не станкостроительный
завод, а запас стали и мощность сталелитейного завода.

\section{Заключение}

В ходе работы задача распределения стали между производством станков,
производством стали и расширением мощностей была сведена к задаче линейного
программирования. Несмотря на поэтапную структуру исходной постановки,
линейность ограничений и переходов состояния позволяет рассматривать весь
пятилетний период как единую LP-модель.

В базовом варианте суммарный выпуск составил 5819.57 станков. Расширение
производственных мощностей в оптимальном плане не используется, так как при
исходных параметрах основной эффект достигается за счет распределения стали
между выпуском станков и производством новой стали.

Исследование начального запаса стали показало, что рост запаса увеличивает
суммарный выпуск, однако после достижения ограничения по мощности
станкостроительного завода дополнительная сталь начинает использоваться для
расширения этой мощности. Исследование начальной мощности станкостроительного
завода показало наличие порогового значения: начиная примерно с 1600 станков
в год дальнейшее увеличение мощности не повышает суммарный выпуск, поскольку
ограничивающим фактором становится ресурсная база и мощность сталелитейного
завода.

Таким образом, метод линейного программирования позволяет не только найти
оптимальный пятилетний план распределения стали, но и исследовать влияние
начальных условий на структуру решения.

\pagebreak
\section{Листинг кода}

\subsection{Реализация на Python}

\lstinputlisting[style=pythonstyle]{lab3_listing.py}


\end{document}
